\section{Uniform variation and the collision-space embeddings}
\label{app:bv-details}
We give the geometric and norm details used in
Section~\ref{sec:collision-covariance}. Variation is always taken in
initial collision coordinates. The statements below do not concern
variation of a density after a many-collision pushforward.

\begin{lemma}[A finite set of possible next centers]
\label{lem:finite-next-centers}
There is a finite subset $\mathcal D$ of the triangular lattice such
that every regular next collision, for every $R\in I$, is with a disk
centered at a member of $\mathcal D$. The cardinality and the set may
be chosen independently of the initial collision state.
\end{lemma}
\begin{proof}
Let $H$ be a uniform upper bound for the free flight, as supplied by the
finite-horizon geometry in Lemma~\ref{lem:collision-spectral-input}.
A flight from $q=Rn_\alpha$ to the disk centered at $d$ ends at
$q'=d+Rn'$ and has length at most $H$. Hence
$|d|\le |q|+|q'-q|+|q'-d|\le H+2R_{\max}$.
The intersection of the lattice with this fixed ball is finite. For
example the open disks of radius $1/2$ about its lattice points are
disjoint, and lie in a ball of radius $H+2R_{\max}+1/2$; area comparison
gives a bound of $4(H+2R_{\max}+1/2)^2$ for their number. Removing the
initial disk or inadmissible centers only decreases this bound.
\end{proof}

\begin{proposition}[Uniform semialgebraic slicing]
\label{prop:semialgebraic-slicing}
Let $\Lambda$ be a semialgebraic parameter set and
$Q=(a,b)\times(c,d)$ a bounded rectangle. Suppose
$u:\Lambda\times Q\to\R$ is semialgebraic and $|u|\le M$.
There is an integer $N$, depending on the family but not on
$\lambda\in\Lambda$, such that $u_\lambda\in\mathrm{BV}(Q)$ and
\begin{equation}\label{eq:slicing-variation-bound}
 |D_1u_\lambda|(Q)\le 4NM(d-c),\qquad
 |D_2u_\lambda|(Q)\le 4NM(b-a).
\end{equation}
The zero extension outside $Q$ has bounded variation as well, with an
additional bound $2M[(b-a)+(d-c)]$ for the sum of its two variations.
The assertions depend only on the almost-everywhere class of the
functions, not on chosen values on lower-dimensional sets.
\end{proposition}
\begin{proof}
Apply differentiable cell decomposition to the graph of $u$, ordering
$(\lambda,y)$ before $x$, and refine the resulting cells by the sign
of the derivative in $x$ wherever that derivative is defined. Uniform
finiteness and one-dimensional monotonicity give a single bound $N$
for the number of monotone intervals and exceptional points in every
fiber. These are the parameterized cell and monotonicity theorems of
\cite[Sections 2.1--2.2 and 6.2]{Coste}. On a monotone interval the
variation is at most $2M$. Jumps between intervals are also at most
$2M$. Isolated assigned values are immaterial for distributional
variation. Thus the variation of the almost-everywhere one-dimensional
fiber is bounded by $4NM$, after enlarging $N$ if necessary.

For $\phi\in C_c^1(Q)$, Fubini's theorem and the one-dimensional
integration-by-parts inequality now give
\[
 \left|\int_Q u_\lambda\,\partial_x\phi\dd x\dd y\right|
 \le 4NM(d-c)\|\phi\|_\infty.
\]
This bounded functional defines a finite distributional derivative
measure, proving the first estimate. Reorder the parameters as
$(\lambda,x)$ before $y$ and repeat the argument for the second.
Taking the larger of the two finite fiber bounds gives the stated
single $N$. Extending each fiber by zero adds at most two endpoint
jumps, each of size at most $M$. The same Fubini argument proves the
claimed zero-extension bound. In particular no bound on a derivative
in the parameter is needed.
\end{proof}

\begin{proposition}[Angular chart gluing for the physical record]
\label{prop:physical-chart-gluing}
The one-collision record $f_R$ has a uniformly bounded supremum norm
and flat-coordinate variation on $M$. Adding any fixed finite union
of moving rectangular section indicators with uniformly bounded
perimeters preserves these bounds. This proves the variation assertions
for $h_R$ and $s_R$ in Lemma~\ref{lem:physical-bv}.
\end{proposition}
\begin{proof}
Choose a fixed finite angular atlas on which
$t=\tan((\alpha-\alpha_0)/2)$ lies in a bounded interval, and choose a
smooth partition of unity with supports strictly inside these charts.
The chart maps and their inverses have bounded first derivatives.
The components of $n_\alpha$ and $t_\alpha$ are rational in $t$, and
$\sqrt{1-p^2}n_\alpha+pt_\alpha$ is semialgebraic in $(t,p)$.
Lemma~\ref{lem:finite-next-centers} reduces the next collision to
finitely many candidate roots. The conditions for a root to be real,
positive, and the first admissible entry, and all comparisons between
two such roots, are semialgebraic. Therefore its minimum and the finite
lattice label define bounded semialgebraic families in $(R,t,p)$.
Ties and singular trajectories can be assigned a fixed bounded value.
They do not change the almost-everywhere record.

Proposition~\ref{prop:semialgebraic-slicing} gives a uniform variation
bound on each chart. Multiplication by a fixed smooth partition
function bounds the variation by its supremum times the original
variation plus its derivative supremum times the original $L^1$ norm.
A bounded smooth change of the angular variable preserves a variation
bound with a constant depending only on the chart and its inverse.
For completeness, this follows for smooth functions by the chain rule
and change of variables, and for a bounded-variation function by
convolution on a slightly smaller chart followed by lower semicontinuity
of distributional variation. The cutoff has support in that smaller
chart. There is therefore no unaccounted jump on the chart boundary.
The sum of the finitely many localized functions gives the asserted
bound on the cylinder. At $p=\pm1$, the fiber proof already includes
finite one-sided traces; no inverse transformation from $p$ to the
angle of incidence is used.

A rectangle indicator has variation equal to its coordinate perimeter,
up to the fixed normalization of $\nu$. The variation of the indicator
of a finite union is bounded by the sum of these perimeters, including
rectangles represented across the angular seam. Its moving center does
not change that estimate. The bounded coefficients in
\eqref{eq:compensation-normalizations} finish the assertion for $h_R$
and $s_R$.
\end{proof}

\begin{lemma}[Smooth densities and test pairings in the local spaces]
\label{lem:density-norm-details}
On each local space $(\mathcal B_j,\mathcal B_{w,j})$ used in
Lemma~\ref{lem:collision-spectral-input}, smooth functions $a,g,b$
satisfy, uniformly for the radii in that neighborhood,
\begin{align}
 \|a\nu\|_{\mathcal B_j}+|a\nu|_{w,j}
       &\le C\|a\|_{C^1},\label{eq:smooth-density-embedding}\\
 \|M_gV\|_{\mathcal B_j}&\le C\|g\|_{C^2}\|V\|_{\mathcal B_j},
 &|\ell_j(M_bV)|&\le C\|b\|_{C^2}\|V\|_{\mathcal B_j}.
 \label{eq:smooth-multiplier-pairing}
\end{align}
The function norms are in the collision coordinates used to define
the local spaces. In particular, for the mean-preserving smoothing,
\begin{equation}\label{eq:smoothed-density-vector}
 \| (\mathsf S_\delta a)\nu\|_{\mathcal B_j}
       \le C\delta^{-2}\|a\|_\infty .
\end{equation}
Here $a$ is a density relative to $\nu$, not relative to unweighted
area in the incidence-angle coordinates.
\end{lemma}
\begin{proof}
We record the norm identification explicitly. In
\cite[Sections 3.2--3.3, (3.11)--(3.13)]{DZ}, a smooth function is
identified with its density times the invariant reference measure
$\mu=c\cos\varphi\dd r\dd\varphi$. Under the local boundary
reparametrization $r=R_0\alpha$, this probability is precisely $\nu$.
Thus the smooth function representing $a\nu$ is $a$; an additional
factor of $\cos\varphi$ is not inserted into that function.

For the weak norm, each stable-curve integral is bounded by
$C\|a\|_\infty$ because the stable curves have uniformly bounded
length and the test functions have bounded supremum. For the strong
stable norm, tests on a curve $W$ have supremum at most $|W|^{-\alpha}$.
The integral is therefore at most
$\|a\|_\infty |W|^{1-\alpha}$, uniformly bounded since $\alpha<1$.
For the strong unstable norm, take two matched graph curves at distance
at most $\epsilon$ in the source metric. Their unmatched parameter
intervals have length $O(\epsilon)$. On their common interval the graph
coordinates and arclength factors differ by $O(\epsilon)$, and the
matched tests differ in supremum by at most $\epsilon$.
The difference of the two integrals is consequently bounded by
$C\|a\|_{C^1}\epsilon$. Dividing by $\epsilon^\beta$ is harmless,
since the source exponent satisfies $\beta<1$. These three estimates
prove \eqref{eq:smooth-density-embedding}, with constants uniform in
the finite local cover and its fixed charts.

For a multiplier, stable and weak tests are simply multiplied by the
restriction of $g$ to their curves. In the matched-curve term, the two
restrictions differ in $C^q$ by at most
$C\|g\|_{C^2}\epsilon^{1-q}$, by interpolation of their supremum
and Lipschitz bounds. The admissible source exponents have
$\beta<1-q$. Splitting into matched tests and this difference gives
the first bound in \eqref{eq:smooth-multiplier-pairing}; it is the
multiplier estimate already used in
Lemma~\ref{lem:collision-spectral-input}. Applying the bounded mass
functional proves the second. Finally
\eqref{eq:bv-smoothing-bounds} bounds the smoothed $C^2$ norm by
$C\delta^{-2}\|a\|_\infty$. Composition with $p=\sin\varphi$
uses only bounded derivatives of this forward map. This proves
\eqref{eq:smoothed-density-vector}.
\end{proof}

The estimates above concern smooth functions on the local collision
spaces. They do not make the hard section projection a bounded
multiplier on those spaces, nor do they realize an unbounded induced
twist. The use of that projection in the compensation argument remains
entirely on the level of ordinary functions and smooth approximations.
